% ============================================================================
% PART II — CONTRIBUTION
% ============================================================================

\part{Contribution}

% ============================================================================
% CHAPTER 3 — PFG: A Graph-Based Framework for Authorization Vulnerability Detection
% ============================================================================

\chapter{PFG: A Graph-Based Framework for Authorization Vulnerability Detection}
\label{ch:theory}

\section{Introduction and Core Insight}
\label{sec:ch3-intro}

Authorization vulnerabilities occupy a unique position in the taxonomy of software security flaws. Unlike injection attacks, which manifest as malicious input propagating through well-defined syntactic channels, or memory corruption, which arises from violations of type safety at the instruction level, authorization vulnerabilities are defined by the \emph{absence} of protective logic along execution paths. A missing \texttt{@login\_required} decorator, an absent ownership check in a queryset filter, or a forgotten role verification before a privileged operation—each represents a structural gap in the security architecture rather than a localized code defect.

This fundamental observation motivates the present work. If authorization vulnerabilities are structural in nature—emerging from the relationship between entry points, protective middleware, and data access operations—then their detection requires a representation that captures this structure explicitly. Token-based approaches that analyze individual functions in isolation cannot reason about cross-component relationships. Generic program representations that model every syntactic detail bury the authorization-relevant signal under irrelevant noise.

This chapter introduces the \emph{Program Flow Graph} (PFG), a domain-specific graph representation designed to capture precisely the architectural elements that participate in authorization decisions. We formalize the graph schema, justify each component of the feature space, develop the mathematical foundations of a hybrid neural-logical detection architecture, and establish the formal rules that encode known vulnerability patterns as structural predicates. The resulting framework provides both the expressiveness of graph neural network classification and the interpretability of rule-based reasoning.

\section{Problem Formalization and Threat Model}
\label{sec:threat-model}

\subsection{Threat Model}

We define the scope of the detection problem through a precise threat model that establishes the attacker capabilities, protected assets, and trust boundaries under consideration.

\paragraph{Attacker Profile.} The adversary is an authenticated user in possession of valid credentials who attempts horizontal or vertical privilege escalation. This distinguishes our threat model from unauthenticated attack scenarios (handled by perimeter security) and from injection-based attacks (handled by input validation). The attacker's goal is to access resources belonging to other users (horizontal escalation), invoke administrative operations without the requisite privilege level (vertical escalation), or access objects across tenant boundaries in multi-tenant deployments.

\paragraph{Protected Assets.} The system under analysis exposes a set of REST API endpoints, each providing programmatic access to application resources. These endpoints accept HTTP requests, process them through a chain of middleware and business logic, and ultimately issue database operations. The protected assets are the data objects accessed by these terminal operations.

\paragraph{Trust Boundary.} The authorization decision chain follows a three-stage model:
\begin{enumerate}
    \item \textbf{Authentication}: verification that the requesting entity possesses a valid identity (session token, JWT, API key).
    \item \textbf{Authorization}: enforcement of access control policies that restrict which authenticated identities may access which resources under which conditions.
    \item \textbf{Data Access}: the terminal database operation that reads, creates, modifies, or deletes a resource.
\end{enumerate}

A vulnerability exists when the authorization stage is absent or insufficient relative to the sensitivity of the data access operation it precedes. The trust boundary lies between authentication and data access; our goal is to detect gaps in this boundary.

\paragraph{Scope.} We target web applications built on the Django REST Framework (DRF) and Spring Boot frameworks—two widely deployed REST API frameworks representing the Python and Java ecosystems respectively. Both implement authorization through well-defined architectural patterns (decorators, permission classes, annotations, security filters) that can be statically extracted, making them amenable to our graph-based analysis.

\subsection{Taxonomy of Authorization Vulnerabilities}

Building on the OWASP Top 10 classification of Broken Access Control (A01:2021) and the Common Weakness Enumeration, we identify five distinct subtypes of authorization vulnerabilities that our framework targets. Each represents a structurally distinct failure mode in the authorization decision chain.

\subsubsection{MISSING\_AUTH (CWE-306)}

An endpoint that performs data access operations without requiring authentication—any anonymous user may invoke the operation and access the underlying data.

\begin{definition}[Missing Authentication]
An endpoint $v$ exhibits MISSING\_AUTH if and only if:
\begin{itemize}
    \item $v$ requires no authentication ($\text{auth}(v) = \text{false}$),
    \item $v$ is reachable to at least one data access node via call edges,
    \item $v$ is not intentionally designated as public (e.g., login, registration).
\end{itemize}
\end{definition}

\textbf{Minimal example:}
\begin{lstlisting}[language=Python, basicstyle=\small\ttfamily]
# urls.py
path('api/documents/<int:pk>/', DocumentView.as_view())
# views.py — no authentication_classes, no permission_classes
class DocumentView(APIView):
    def get(self, request, pk):
        return Response(Document.objects.get(pk=pk))
\end{lstlisting}

\subsubsection{MISSING\_ROLE (CWE-862)}

An authenticated endpoint that performs privileged operations (mutations on shared resources) without verifying that the requesting user holds the requisite privilege level.

\begin{definition}[Missing Role Check]
An endpoint $v$ exhibits MISSING\_ROLE if:
\begin{itemize}
    \item $v$ requires authentication ($\text{auth}(v) = \text{true}$),
    \item $v$ reaches a data access node performing a privileged operation (WRITE, UPDATE, DELETE) on a shared entity,
    \item No node on the call path from $v$ to the data access enforces a role check with minimum privilege USER or above.
\end{itemize}
\end{definition}

\subsubsection{IDOR — Insecure Direct Object Reference (CWE-639)}

An endpoint accepts a user-controlled identifier (route parameter) that flows to a database query without verifying that the authenticated user owns or has permission to access the referenced object.

\begin{definition}[Insecure Direct Object Reference]
An endpoint $v$ exhibits IDOR if:
\begin{itemize}
    \item $v$ contains route parameters (e.g., \texttt{<int:pk>}),
    \item An identifier flow edge connects $v$ to a data access node $d$,
    \item The data access operation targets an ownership-bearing entity,
    \item No node on any call path from $v$ to $d$ enforces an ownership check,
    \item The data access node has no ownership predicate in its query filter.
\end{itemize}
\end{definition}

\textbf{Minimal example:}
\begin{lstlisting}[language=Python, basicstyle=\small\ttfamily]
class WorkoutDetailView(RetrieveUpdateDestroyAPIView):
    permission_classes = [IsAuthenticated]  # authenticated but...
    def get_queryset(self):
        return Workout.objects.all()  # no .filter(user=request.user)
\end{lstlisting}

\subsubsection{TENANT\_ISOLATION\_FAILURE}

In multi-tenant applications, a data access operation retrieves resources across tenant boundaries without enforcing tenant-scoped filtering.

\begin{definition}[Tenant Isolation Failure]
An endpoint $v$ exhibits TENANT\_ISOLATION\_FAILURE if:
\begin{itemize}
    \item $v$ reaches a data access node $d$ whose target entity requires tenant isolation,
    \item $d$ has no tenant predicate in its query filter,
    \item No middleware on the path from $v$ to $d$ enforces tenant scoping.
\end{itemize}
\end{definition}

\subsubsection{OBJECT\_LEVEL\_PERMISSION\_BYPASS}

An endpoint has model-level permission evidence (the application defines object-level permissions for the resource) but the specific endpoint does not enforce object-level authorization when accessing individual objects via route parameters.

\begin{definition}[Object-Level Permission Bypass]
An endpoint $v$ exhibits OBJECT\_LEVEL\_PERMISSION\_BYPASS if:
\begin{itemize}
    \item $v$ has model-level permission indicators ($\text{model\_perm}(v) = \text{true}$),
    \item $v$ accepts route parameters and reaches a data access via identifier flow,
    \item No object-level permission check is enforced on any call path to the data access.
\end{itemize}
\end{definition}

\subsection{Research Questions}

Given the threat model and vulnerability taxonomy established above, this thesis addresses the following research questions:

\begin{description}
    \item[RQ1:] Can a domain-specific graph representation effectively detect access control vulnerabilities in production web applications?
    \item[RQ2:] Does the framework provide interpretable, actionable results without requiring post-hoc explanation techniques?
    \item[RQ3:] How does GNN-based detection compare with rule-based structural analysis and existing static analysis security tools?
\end{description}

\section{Positioning Against Existing Paradigms}
\label{sec:positioning}

Before presenting our approach, we establish precisely \emph{why} existing detection paradigms are insufficient for authorization vulnerabilities. This argumentation is made once here and will not be repeated in subsequent chapters.

\subsection{Against NLP-Based Detection}

Recent advances in pre-trained language models for code—CodeBERT \cite{feng2020codebert}, VulBERTa \cite{hanif2022vulberta}, LineVul \cite{fu2022linevul}—have demonstrated strong performance on vulnerability detection benchmarks, particularly for injection and memory safety vulnerabilities. These models operate by encoding source code as token sequences within fixed-length windows (typically 512 tokens) and learning vulnerability-correlated patterns in the token distribution.

Authorization vulnerabilities are fundamentally incompatible with this paradigm for three reasons:

\begin{enumerate}
    \item \textbf{Cross-file dependency.} An authorization decision involves at minimum three files: the URL configuration (\texttt{urls.py}), the view handler (\texttt{views.py}), and the data model (\texttt{models.py}). The security-critical relationship between these components—whether an endpoint's route correctly maps through a permission-checking middleware to a properly-filtered data access—cannot be captured within a single-function token window.

    \item \textbf{Absence-based signal.} Token-based models learn to detect the \emph{presence} of vulnerability-indicative patterns. Authorization flaws are defined by what is \emph{missing}: a decorator that should be present, a filter clause that should constrain a queryset, a permission class that should guard an endpoint. Learning the absence of a pattern within a space of possible patterns is fundamentally harder than learning presence.

    \item \textbf{Structural, not lexical, signal.} The vulnerability signal in an IDOR lies not in any particular token or token sequence, but in the \emph{structural relationship} between a route parameter, an identifier flow path, and the absence of an ownership predicate at the data access terminus. This relationship is topological, not sequential.
\end{enumerate}

\subsection{Against Generic Code Property Graphs}

Code Property Graphs (CPGs), as implemented in tools like Joern \cite{yamaguchi2014modeling}, compose the Abstract Syntax Tree, Control Flow Graph, and Program Dependence Graph into a unified representation. CPG-based detection has proven effective for memory safety vulnerabilities in C/C++ where taint flows from source to sink can be expressed as graph reachability queries.

For authorization vulnerability detection, CPGs suffer from a fundamental signal-to-noise problem:

\begin{enumerate}
    \item \textbf{Representation granularity mismatch.} A CPG represents every syntactic element, control flow edge, and data dependency at the statement level. For a typical Django view function of 20 lines, the CPG contains several hundred nodes. Yet the authorization-relevant information can be expressed in 3--5 nodes: one endpoint node, zero or more middleware guards, and one or more data access operations. The security signal—``does a permission check exist between this entry point and this database query''—is buried under 99\% irrelevant structure.

    \item \textbf{Framework-agnostic representation.} CPGs model language-level semantics (function calls, assignments, branches) but not framework-level abstractions (URL routing, decorator-based permission enforcement, queryset filtering). The architectural patterns that constitute authorization in Django or Spring cannot be queried without framework-specific semantic overlays that transform the CPG into something closer to our purpose-built abstraction.

    \item \textbf{Scalability.} CPG construction for a large web application (500+ endpoints) produces graphs with millions of nodes. GNN inference over such graphs is computationally prohibitive and yields diminishing returns because the vast majority of the graph structure is irrelevant to the authorization detection task.
\end{enumerate}

\subsection{Our Positioning}

The PFG occupies a distinct position in the design space:

\begin{itemize}
    \item \textbf{Purpose-built abstraction}: Rather than modeling all program structure and filtering post-hoc, the PFG models \emph{only} the four node types that participate in authorization decisions. A typical PFG contains 5--50 nodes, each representing a meaningful architectural component.

    \item \textbf{Interpretable by construction}: Every node feature is a human-readable boolean or categorical security property (e.g., \texttt{requires\_auth}, \texttt{ownership\_check\_present}). No opaque learned embeddings are required—the feature space is fully transparent.

    \item \textbf{Structural explanation without post-hoc tools}: Because each node maps to an architectural component and each edge represents an explicit relationship (call, identifier flow, tenant flow), a detection can be explained as a concrete path: ``Endpoint $E$ at route \texttt{/api/items/<pk>/} has identifier flow to data access $D$ on entity \texttt{Item} with no ownership guard on any intervening path.'' No GNNExplainer \cite{ying2019gnnexplainer} or other post-hoc technique is needed.
\end{itemize}

Both opaque embeddings and post-hoc explainability are eliminated \emph{by design}, not compensated for after the fact.

\section{The Program Flow Graph Formulation}
\label{sec:pfg-formulation}

\subsection{Formal Graph Definition}

\begin{definition}[Program Flow Graph]
A Program Flow Graph is a directed, typed, attributed graph defined as the 6-tuple:
\begin{equation}
G = (V, E, \tau_V, \tau_E, \phi_V, \phi_E)
\end{equation}
where:
\begin{itemize}
    \item $V$ is the finite set of nodes,
    \item $E \subseteq V \times V$ is the set of directed edges,
    \item $\tau_V: V \to \{\textsc{endpoint}, \textsc{middleware}, \textsc{service}, \textsc{data\_access}\}$ assigns each node its type,
    \item $\tau_E: E \to \{\textsc{call}, \textsc{identifier\_flow}, \textsc{tenant\_flow}\}$ assigns each edge its type,
    \item $\phi_V: V \to \mathbb{R}^{36}$ maps each node to its feature vector,
    \item $\phi_E: E \to \mathbb{R}^{3}$ maps each edge to its feature vector (one-hot encoding of edge type).
\end{itemize}
\end{definition}

The PFG schema is formally specified as a JSON Schema (version 1.0) and all generated graphs are validated against this schema at construction time, ensuring structural consistency across the pipeline.

\subsection{Node Types and Their Semantic Role}

Each node type corresponds to a specific role in the authorization decision chain. Together, they partition the application's security-relevant architecture into four orthogonal concerns.

\subsubsection{ENDPOINT Nodes}

An \textsc{endpoint} node represents the entry point where an HTTP request enters the application. It encodes the route pattern, accepted HTTP methods, route parameters, and initial authorization signals derived from the framework's configuration (decorators, permission classes, mixin inheritance).

Formally, for an endpoint node $v$ with $\tau_V(v) = \textsc{endpoint}$:
\begin{equation}
\text{attrs}(v) = \langle \text{method}, \text{route}, \text{params}, \text{auth}, \text{role}, \text{ownership}, \ldots \rangle
\end{equation}

The endpoint is the origin of all authorization reasoning: \emph{what protections should this entry point have, given what data it ultimately accesses?}

\subsubsection{MIDDLEWARE Nodes}

A \textsc{middleware} node represents a protective guard interposed between the entry point and the data access. Each middleware has a \emph{guard type} drawn from five categories:
\begin{itemize}
    \item \textsc{authentication}: verifies the requester's identity,
    \item \textsc{role\_check}: verifies the requester holds a minimum privilege level,
    \item \textsc{ownership\_check}: verifies the requester owns the target object,
    \item \textsc{tenant\_check}: enforces tenant boundary isolation,
    \item \textsc{object\_permission\_check}: enforces fine-grained object-level permissions.
\end{itemize}

The presence or absence of middleware nodes on paths from endpoints to data access nodes constitutes the primary vulnerability signal in the PFG.

\subsubsection{SERVICE Nodes}

A \textsc{service} node represents intermediate business logic invoked along the call chain between an endpoint and its data access operations. Service nodes model the functional layers of the application (business services, utility functions, nested view calls) and serve as routing vertices in the call graph. While not directly involved in authorization decisions, they enable accurate path reconstruction for vulnerability explanation.

\subsubsection{DATA\_ACCESS Nodes}

A \textsc{data\_access} node represents a terminal database operation that reads, writes, updates, or deletes a resource. Its attributes encode:
\begin{itemize}
    \item The \emph{operation type}: \textsc{read}, \textsc{write}, \textsc{update}, or \textsc{delete},
    \item The \emph{target entity}: the database model being accessed,
    \item \emph{Ownership predicates}: whether the query filters by the requesting user,
    \item \emph{Tenant predicates}: whether the query scopes to the requesting tenant,
    \item \emph{Object permission indicators}: whether object-level checks are applied.
\end{itemize}

\paragraph{Vulnerability as a path property.} With these four node types defined, we can state the central insight of the PFG formalization:

\begin{quote}
\emph{A vulnerability exists when a path from an \textsc{endpoint} to a \textsc{data\_access} node lacks sufficient \textsc{middleware} given the sensitivity of the data access operation.}
\end{quote}

\subsection{Edge Types}

Three edge types capture the relationships between nodes:

\begin{itemize}
    \item \textbf{CALL} edges ($\tau_E = \textsc{call}$): Represent function invocation, connecting endpoints to services, services to other services, and services to data access nodes. These edges form the call graph backbone of the PFG.

    \item \textbf{IDENTIFIER\_FLOW} edges ($\tau_E = \textsc{identifier\_flow}$): Track how user-controlled identifiers (route parameters) propagate from endpoints to data access operations. The presence of an identifier flow edge from an endpoint to a data access node is a necessary condition for an IDOR vulnerability.

    \item \textbf{TENANT\_FLOW} edges ($\tau_E = \textsc{tenant\_flow}$): Track tenant-scoping information flow from authentication context to data access queries. Their absence on paths to tenant-sensitive data indicates a potential tenant isolation failure.
\end{itemize}

\subsection{Feature Vector Justification}

The node feature function $\phi_V: V \to \mathbb{R}^{36}$ encodes each node as a 36-dimensional vector. Unlike learned embeddings in NLP-based approaches, every dimension corresponds to a human-interpretable security property. This design choice ensures that the feature space is fully transparent and auditable.

Table~\ref{tab:feature-vector} provides the complete specification of the feature vector, mapping each index to its semantic meaning and the vulnerability subtype it is designed to detect.

\begin{table}[htbp]
\centering
\caption{PFG node feature vector specification ($\phi_V \in \mathbb{R}^{36}$).}
\label{tab:feature-vector}
\small
\begin{tabular}{@{}clll@{}}
\toprule
\textbf{Index} & \textbf{Property} & \textbf{Encoding} & \textbf{Detection Role} \\
\midrule
$[0{:}4]$ & Node type & One-hot & Graph structure \\
$[4{:}9]$ & HTTP method & One-hot & Risk profile \\
$[9]$ & \texttt{requires\_auth} & Binary & MISSING\_AUTH \\
$[10{:}15]$ & Required role & One-hot & MISSING\_ROLE \\
$[15]$ & Ownership check & Binary & IDOR \\
$[16{:}19]$ & Guard type & One-hot & Path protection \\
$[19{:}23]$ & Operation type & One-hot & Damage potential \\
$[23]$ & Ownership predicate & Binary & IDOR (data-level) \\
$[24]$ & Has route params & Binary & IDOR prerequisite \\
$[25]$ & Reachable data count & Normalized & Exposure scope \\
$[26]$ & Has identifier flow & Binary & IDOR prerequisite \\
$[27]$ & Unprotected data & Binary & Multi-subtype \\
$[28]$ & Role check present & Binary & MISSING\_ROLE \\
$[29]$ & Model permission & Binary & OBJ\_PERM\_BYPASS \\
$[30]$ & Object permission & Binary & OBJ\_PERM\_BYPASS \\
$[31]$ & Tenant check & Binary & TENANT\_ISO \\
$[32]$ & Has tenant scope & Binary & TENANT\_ISO \\
$[33]$ & Unscoped tenant data & Binary & TENANT\_ISO \\
$[34]$ & Obj-perm data reach & Binary & OBJ\_PERM\_BYPASS \\
$[35]$ & Tenant/obj signal & Continuous & Multi-subtype \\
\bottomrule
\end{tabular}
\end{table}

\paragraph{Design rationale.} The feature vector is structured in layers of increasing specificity:
\begin{itemize}
    \item \textbf{Structural identity} (indices 0--4): node type and HTTP method establish the node's role in the graph and the risk profile of the operation (DELETE/PUT operations are higher-risk than GET).
    \item \textbf{Direct authorization signals} (indices 9--15): properties directly read from the framework configuration that indicate whether an authorization mechanism is configured.
    \item \textbf{Guard evidence} (indices 16--18): for middleware nodes, the type of protection they provide.
    \item \textbf{Data access characterization} (indices 19--23): the nature and sensitivity of the terminal operation.
    \item \textbf{Structural reachability signals} (indices 24--35): computed post-construction features that encode relationships between the node and its graph neighborhood—whether unprotected data is reachable, whether tenant boundaries are respected, whether identifier flows exist without corresponding guards.
\end{itemize}

The critical property of this feature space is that \emph{every dimension has a semantically meaningful interpretation} in the context of authorization security. This enables not only GNN-based classification but also human audit of the features contributing to any given prediction.

\section{Hybrid Detection Architecture}
\label{sec:architecture}

The PFG framework employs a hybrid architecture combining a graph neural network for vulnerability scoring with a structural rule engine for validation and explanation. This section presents the mathematical foundations of both components and establishes their complementarity.

\subsection{GATv2-Based Classification}

We select Graph Attention Network v2 (GATv2) \cite{brody2021attentive} as the message-passing backbone for the neural component. The choice is motivated by the asymmetric nature of authorization relationships in the PFG: the protective effect of a middleware node on an endpoint depends on \emph{which} data access the endpoint reaches, and the severity of an unprotected path depends on the sensitivity of the terminal operation. Static attention coefficients (as in the original GAT \cite{velivckovic2018graph}) cannot capture this context dependence; GATv2's dynamic attention mechanism computes attention as a function of both source and target node features.

\subsubsection{Message-Passing Equations}

Given a graph $G = (V, E)$ with node features $\mathbf{h}_i^{(l)} \in \mathbb{R}^d$ at layer $l$ and edge features $\mathbf{e}_{ij} \in \mathbb{R}^3$, the GATv2 layer computes updated representations through the following attention mechanism:

\begin{equation}
\label{eq:gatv2-attention}
\alpha_{ij}^{(l)} = \frac{\exp\left(\mathbf{a}^{\top} \text{LeakyReLU}\left(\mathbf{W}\left[\mathbf{h}_i^{(l)} \| \mathbf{h}_j^{(l)} \| \mathbf{e}_{ij}\right]\right)\right)}{\sum_{k \in \mathcal{N}(i)} \exp\left(\mathbf{a}^{\top} \text{LeakyReLU}\left(\mathbf{W}\left[\mathbf{h}_i^{(l)} \| \mathbf{h}_k^{(l)} \| \mathbf{e}_{ik}\right]\right)\right)}
\end{equation}

where $\mathbf{W} \in \mathbb{R}^{d' \times (2d + 3)}$ is the shared weight matrix, $\mathbf{a} \in \mathbb{R}^{d'}$ is the attention vector, $\|$ denotes concatenation, and $\mathcal{N}(i)$ is the neighborhood of node $i$.

The updated node representation aggregates neighbor information weighted by the computed attention:

\begin{equation}
\label{eq:gatv2-update}
\mathbf{h}_i^{(l+1)} = \Big\|_{k=1}^{K} \sum_{j \in \mathcal{N}(i)} \alpha_{ij}^{(l,k)} \mathbf{W}^{(k)} \mathbf{h}_j^{(l)}
\end{equation}

where $K$ is the number of attention heads and $\|$ denotes concatenation across heads.

\paragraph{Why attention matters for authorization.} In the PFG, the protective relationship between nodes is inherently asymmetric and context-dependent. Consider a middleware node $m$ of type \textsc{ownership\_check} on a path between endpoint $v$ and data access $d$. The relevance of $m$ to the vulnerability classification of $v$ depends on:
\begin{enumerate}
    \item Whether $d$ is an ownership-bearing entity (if not, the ownership check is irrelevant),
    \item Whether the identifier flow actually passes through $m$ (if the flow bypasses the guard, it provides no protection),
    \item Whether other guards on the same path provide equivalent protection.
\end{enumerate}

The dynamic attention mechanism in GATv2 allows the model to learn these context-dependent relevance weights, assigning high attention to guards that are actually effective for the specific vulnerability subtype being evaluated and low attention to structurally present but semantically irrelevant protections.

\subsubsection{Architecture Configuration}

The PathAnalyzer model instantiates the GATv2 backbone with the following configuration:
\begin{itemize}
    \item Input projection: $\text{Linear}(36 \to 64)$
    \item 3 GATv2 layers, each with 4 attention heads and per-head dimension $64/4 = 16$
    \item Concatenation across heads restores dimension 64
    \item ELU activation with dropout $p = 0.1$ after each layer
    \item Edge features ($\mathbb{R}^3$) incorporated into the attention computation
\end{itemize}

\subsection{Target-Aware Attention Pooling}

After $L$ layers of message passing, each node $i$ has a contextualized representation $\mathbf{h}_i^{(L)} \in \mathbb{R}^{64}$ that encodes both its own features and information propagated from its neighborhood. To classify a specific endpoint, we must aggregate the entire graph's information into a fixed-size representation while maintaining focus on the target endpoint being evaluated.

Standard global pooling approaches (mean, max, or sum over all nodes) lose this target-specificity: the same pooled representation would be used regardless of which endpoint is being classified. This is problematic because a graph containing multiple endpoints may be simultaneously safe for some and vulnerable for others.

We introduce a \emph{target-aware attention pooling} mechanism that computes a graph-level representation anchored to the specific target endpoint:

\begin{definition}[Target-Aware Attention Pooling]
Given node representations $\{\mathbf{h}_i\}_{i \in V}$ and target endpoint representation $\mathbf{h}_{\text{target}}$ (mean of representations at target-masked nodes), the pooled graph representation is:
\end{definition}

\begin{equation}
\label{eq:attn-score}
s_i = \mathbf{w}^{\top} \tanh\left(\mathbf{W}_s \left[\mathbf{h}_i \| \mathbf{h}_{\text{target}}\right]\right)
\end{equation}

\begin{equation}
\label{eq:attn-weight}
\alpha_i = \frac{\exp(s_i)}{\sum_{j \in V} \exp(s_j)}
\end{equation}

\begin{equation}
\label{eq:pooled}
\mathbf{z} = \sum_{i \in V} \alpha_i \, \mathbf{h}_i
\end{equation}

\begin{equation}
\label{eq:classification}
\hat{y} = \text{MLP}\left(\left[\mathbf{z} \| \mathbf{h}_{\text{target}}\right]\right)
\end{equation}

where $\mathbf{W}_s \in \mathbb{R}^{64 \times 128}$ and $\mathbf{w} \in \mathbb{R}^{64}$ are learnable parameters, and the final MLP maps from $\mathbb{R}^{128}$ to $\mathbb{R}^{C}$ where $C = 8$ is the number of output classes.

\paragraph{Interpretation.} The attention score $s_i$ measures how relevant each node $i$ is to the vulnerability classification of the target endpoint. Nodes that are structurally important to the target's security posture (its guards, its data access operations, intermediate services on its paths) receive higher attention weights, while nodes belonging to other subgraphs in the same application are down-weighted. This mechanism ensures that the model's classification focuses on the local security context of the specific endpoint being evaluated.

\paragraph{Output classes.} The classifier produces scores over 8 classes: \textsc{safe}, \textsc{missing\_auth}, \textsc{missing\_role}, \textsc{idor}, \textsc{tenant\_isolation\_failure}, \textsc{object\_level\_permission\_bypass}, \textsc{combined}, and \textsc{intentionally\_public}. The \textsc{combined} class handles endpoints exhibiting multiple simultaneous vulnerabilities.

\subsection{Five Structural Rules as Formal Boolean Predicates}

The rule-based component provides guaranteed interpretability and serves as both a standalone baseline and a validation layer for GNN predictions. Each rule is formalized as a first-order logic predicate over the PFG structure.

\begin{definition}[Rule $R_1$ — Missing Authentication]
\begin{equation}
R_1: \exists v \in V \;\Big|\; \tau_V(v) = \textsc{ep} \wedge \neg\text{auth}(v) \wedge \text{reaches}(v, \textsc{da}) \wedge \neg\text{public}(v) \implies \textsc{missing\_auth}
\end{equation}
\end{definition}

\begin{definition}[Rule $R_2$ — Missing Role Check]
\begin{equation}
R_2: \exists v \;\Big|\; \text{auth}(v) \wedge \text{reaches\_privileged}(v) \wedge \neg\exists m \in \text{path}(v, d) : \text{guard}(m) = \textsc{role} \implies \textsc{missing\_role}
\end{equation}
\end{definition}

\begin{definition}[Rule $R_3$ — IDOR]
\begin{equation}
R_3: \exists v \;\Big|\; \text{has\_param}(v) \wedge \exists (v, d) \in E_{\textsc{id\_flow}} \wedge \neg\text{ownership\_on\_path}(v, d) \implies \textsc{idor}
\end{equation}
\end{definition}

\begin{definition}[Rule $R_4$ — Tenant Isolation Failure]
\begin{equation}
R_4: \exists v \;\Big|\; \text{reaches\_tenant\_data}(v) \wedge \neg\text{tenant\_scoped}(v, d) \implies \textsc{tenant\_iso}
\end{equation}
\end{definition}

\begin{definition}[Rule $R_5$ — Object-Level Permission Bypass]
\begin{equation}
R_5: \exists v \;\Big|\; \text{model\_perm}(v) \wedge \text{has\_param}(v) \wedge \neg\text{obj\_perm\_on\_path}(v, d) \implies \textsc{obj\_perm\_bypass}
\end{equation}
\end{definition}

where the predicates are defined as:
\begin{itemize}
    \item $\text{reaches}(v, \textsc{da})$: $\exists d \in V, \tau_V(d) = \textsc{data\_access}$ reachable from $v$ via call edges,
    \item $\text{reaches\_privileged}(v)$: $v$ reaches a data access performing WRITE, UPDATE, or DELETE,
    \item $\text{path}(v, d)$: the set of all simple paths from $v$ to $d$ via call edges (bounded by depth 8),
    \item $\text{ownership\_on\_path}(v, d)$: $\exists m \in \text{path}(v, d)$ such that $m$ is an ownership guard or $d$ has an ownership predicate,
    \item $\text{tenant\_scoped}(v, d)$: tenant check on path or tenant predicate at $d$.
\end{itemize}

\subsection{Complementarity of Neural and Logical Components}

The two components serve distinct and complementary roles:

\begin{itemize}
    \item \textbf{Rule baseline — precision and interpretability}: The five rules are complete for explicitly encodable structural patterns. Every rule firing produces a natural-language explanation grounded in the specific nodes and edges involved. Rules provide guaranteed zero false positives within their scope: if the structural conditions are satisfied, the vulnerability is by definition present.

    \item \textbf{GNN — generalization and recall}: The neural component captures interaction effects and contextual patterns that exceed the expressiveness of boolean predicates. For instance, the GNN may learn that certain combinations of features across a multi-hop neighborhood indicate a vulnerability even when no single rule fires—such as a chain of services with decreasing privilege levels terminating in an unguarded data access.

    \item \textbf{Combined pipeline}: In the integrated system, the GNN proposes candidate vulnerabilities (providing recall), and the structural validator confirms predictions against program invariants (providing precision and explanation). Only predictions that satisfy at least one structural rule are reported to the user, ensuring every output is accompanied by a concrete, verifiable explanation.
\end{itemize}

\section{Interpretability by Design}
\label{sec:interpretability}

A fundamental requirement for security tooling is that developers can understand \emph{why} a particular finding was raised. Tools that produce opaque vulnerability scores without actionable explanations are rarely adopted in practice, regardless of their accuracy.

The PFG framework achieves interpretability through three mechanisms:

\begin{enumerate}
    \item \textbf{Transparent features}: Every dimension of the feature vector corresponds to a named security property. A developer inspecting a prediction can immediately identify which features contributed (e.g., ``this endpoint has \texttt{requires\_auth=False} and \texttt{reachable\_data\_count=3}'').

    \item \textbf{Structural path explanation}: Each \texttt{Finding} produced by the validator traces an explicit path through the PFG: a specific endpoint, the absence of a specific guard type, and the specific data access operation at risk. The explanation is derived from graph structure, not from learned attention weights.

    \item \textbf{Actionable remediation}: Because each finding identifies what is missing (which guard type, on which path, protecting which data), the remediation is immediate: add the specified guard. For IDOR findings, the fix is to add an ownership predicate to the queryset. For MISSING\_AUTH, the fix is to add an authentication class.
\end{enumerate}

\paragraph{Contrast with post-hoc explainability.} Approaches based on GNNExplainer \cite{ying2019gnnexplainer} or attention visualization learn \emph{which subgraph was most important} for a prediction after the fact. These explanations are approximate, may not correspond to the actual causal mechanism, and require additional computational overhead. Our structural validator \emph{constructs} the explanation from the program invariants that the vulnerability violates—the explanation IS the detection mechanism, not a separate approximation.

This design satisfies RQ2 by construction: interpretability is a first-class property of the architecture, not an afterthought.

\section{Conclusion}
\label{sec:ch3-conclusion}

This chapter has established the theoretical foundations of the PFG framework. We formalized the threat model targeting authenticated privilege escalation in REST APIs, defined a vulnerability taxonomy of five structurally distinct authorization flaw subtypes, and demonstrated why existing detection paradigms (NLP-based and CPG-based) are fundamentally mismatched to this problem class.

The Program Flow Graph was defined as a 6-tuple $(V, E, \tau_V, \tau_E, \phi_V, \phi_E)$ with four node types, three edge types, a 36-dimensional interpretable feature vector, and three edge features. Upon this representation, we developed a hybrid detection architecture combining GATv2 message-passing with target-aware attention pooling for neural classification, and five structural rules formalized as first-order logic predicates for validation and explanation.

The next chapter grounds this theoretical framework in engineering reality: we describe how PFGs are extracted from real Django and Spring Boot applications, how the evaluation corpus of 75 real-world CVEs was constructed, and how data augmentation and training methodology address the extreme class imbalance inherent to vulnerability detection.


% ============================================================================
% CHAPTER 4 — Implementation and Dataset
% ============================================================================

\chapter{Implementation and Dataset}
\label{ch:implementation}

\section{Introduction}
\label{sec:ch4-intro}

The theoretical framework presented in Chapter~\ref{ch:theory} defines \emph{what} a Program Flow Graph represents and \emph{why} its components are chosen. This chapter addresses the engineering challenge of \emph{how}: how PFGs are extracted from real application source code, how the evaluation corpus was constructed from production open-source software, how class imbalance is addressed through data augmentation, and how the training pipeline is configured. We also briefly present the demonstration system that provides real-time visualization of the analysis process.

Table~\ref{tab:tech-stack} summarizes the key technologies employed in the implementation.

\begin{table}[htbp]
\centering
\caption{Technology stack.}
\label{tab:tech-stack}
\small
\begin{tabular}{@{}lll@{}}
\toprule
\textbf{Component} & \textbf{Technology} & \textbf{Justification} \\
\midrule
Static extraction & Python \texttt{ast} module & Framework-native AST access \\
Spring extraction & Regex-based parsing & Annotation-centric patterns \\
Schema validation & Pydantic v2 & Runtime type enforcement \\
GNN training & PyTorch Geometric & Production GNN framework \\
API backend & FastAPI & Async WebSocket support \\
Frontend & Next.js 14 & Server-side rendering, React \\
Graph rendering & D3.js force layout & Interactive graph visualization \\
\bottomrule
\end{tabular}
\end{table}

\section{Framework-Specific PFG Extraction}
\label{sec:extractors}

PFG extraction transforms raw application source code into the formal graph structure defined in Chapter~\ref{ch:theory}. This process is inherently framework-specific: Django and Spring Boot implement authorization through fundamentally different architectural patterns, requiring distinct extraction strategies.

\subsection{Django Extractor (AST-Based)}

The Django extractor operates on the Python Abstract Syntax Tree and comprises five sub-modules, each addressing a specific aspect of PFG construction.

\subsubsection{Endpoint Discovery}

Django applications define URL-to-view mappings through a tree-structured URL configuration system. The root \texttt{urls.py} module contains \texttt{urlpatterns}—a list of \texttt{path()}, \texttt{re\_path()}, and \texttt{include()} directives that may nest arbitrarily deep through multiple modules.

The endpoint discovery algorithm performs a recursive traversal of this tree, resolving:
\begin{itemize}
    \item Direct \texttt{path()} and \texttt{re\_path()} entries mapping routes to view functions or classes,
    \item \texttt{include()} directives that delegate to sub-modules with prefix accumulation,
    \item DRF \texttt{Router} instances that generate multiple URL patterns from a single \texttt{ViewSet} registration,
    \item Class-based view method resolution (\texttt{as\_view()} dispatch to \texttt{get()}, \texttt{post()}, etc.),
    \item Custom \texttt{@action} decorators on ViewSets that generate additional routes.
\end{itemize}

\begin{algorithm}
\caption{Endpoint Discovery via URLConf Traversal}
\label{alg:endpoint-discovery}
\begin{algorithmic}[1]
\REQUIRE Root URLConf module path $M$, accumulated prefix $P = $ ``''
\ENSURE Set of endpoint descriptors $\mathcal{E}$
\STATE $\mathcal{E} \leftarrow \emptyset$
\STATE $\text{patterns} \leftarrow \text{parse\_urlpatterns}(M)$
\FOR{each pattern $p$ in patterns}
    \IF{$p$ is \texttt{include(submodule, prefix)}}
        \STATE $\mathcal{E} \leftarrow \mathcal{E} \cup \text{DiscoverEndpoints}(\text{submodule}, P + p.\text{route})$
    \ELSIF{$p$ is \texttt{Router} registration}
        \STATE $\text{viewset} \leftarrow \text{resolve\_viewset}(p)$
        \FOR{each action $a$ in viewset.actions $\cup$ viewset.custom\_actions}
            \STATE $\mathcal{E} \leftarrow \mathcal{E} \cup \{(P + p.\text{prefix} + a.\text{suffix}, a.\text{method}, \text{viewset})\}$
        \ENDFOR
    \ELSE
        \STATE $\text{view} \leftarrow \text{resolve\_view}(p.\text{callback})$
        \STATE $\mathcal{E} \leftarrow \mathcal{E} \cup \{(P + p.\text{route}, p.\text{methods}, \text{view})\}$
    \ENDIF
\ENDFOR
\RETURN $\mathcal{E}$
\end{algorithmic}
\end{algorithm}

\subsubsection{Authorization Annotation}

For each discovered endpoint, the extractor determines its authorization configuration by inspecting multiple framework-specific mechanisms:

\begin{itemize}
    \item \textbf{Decorator inspection}: \texttt{@login\_required}, \texttt{@permission\_required('perm')}, \texttt{@user\_passes\_test(fn)} on function-based views.
    \item \textbf{Permission classes}: \texttt{permission\_classes = [IsAuthenticated, IsAdminUser]} on class-based DRF views.
    \item \textbf{Mixin inheritance}: \texttt{LoginRequiredMixin}, \texttt{PermissionRequiredMixin} in the class hierarchy.
    \item \textbf{Settings defaults}: \texttt{REST\_FRAMEWORK['DEFAULT\_PERMISSION\_CLASSES']} and \texttt{DEFAULT\_AUTHENTICATION\_CLASSES} in project settings.
    \item \textbf{ViewSet-level overrides}: Per-action \texttt{permission\_classes} via \texttt{get\_permissions()} method or \texttt{@action(permission\_classes=[...])} decorators.
\end{itemize}

The authorization annotation produces the initial values for the endpoint's feature vector: \texttt{requires\_auth}, \texttt{required\_role}, \texttt{ownership\_check\_present}, and related properties.

\subsubsection{Call Graph Construction}

The call graph connects endpoints to the services and data access operations they invoke. Construction proceeds through intra-project AST traversal with the following algorithm:

\begin{enumerate}
    \item For each endpoint's handler function/method, parse its AST and extract all call expressions.
    \item Resolve each call target to a module-qualified function name using import alias resolution.
    \item If the target is within the project scope, recursively process its body.
    \item Continue until reaching a maximum depth of 8 or encountering only external library calls.
\end{enumerate}

Import alias resolution handles the common Django patterns of renamed imports (\texttt{from myapp.services import UserService as US}) and lazy references (\texttt{apps.get\_model('myapp', 'MyModel')}).

\subsubsection{Data Access Detection}

Data access nodes are created when the call graph traversal encounters Django ORM operations:
\begin{itemize}
    \item \texttt{Model.objects.get()}, \texttt{.filter()}, \texttt{.all()}, \texttt{.exclude()} $\to$ \textsc{read}
    \item \texttt{Model.objects.create()}, \texttt{Model.save()} $\to$ \textsc{write}
    \item \texttt{.update()}, \texttt{Model.save()} on existing $\to$ \textsc{update}
    \item \texttt{.delete()} $\to$ \textsc{delete}
\end{itemize}

For each detected data access, the extractor classifies its ownership predicate into three categories:
\begin{itemize}
    \item \textbf{DIRECT}: The query explicitly filters by the request user (e.g., \texttt{.filter(user=request.user)}),
    \item \textbf{INDIRECT}: The query filters by a related field that implies ownership (e.g., \texttt{.filter(author\_\_org=request.user.org)}),
    \item \textbf{NONE}: No ownership constraint is present.
\end{itemize}

\subsubsection{Identifier Flow Tracking}

Identifier flow edges track how route parameters propagate from URL patterns to database queries. The algorithm performs intra-procedural taint propagation:

\begin{algorithm}
\caption{Identifier Flow Taint Propagation}
\label{alg:id-flow}
\begin{algorithmic}[1]
\REQUIRE Endpoint $v$ with route parameters $P = \{p_1, \ldots, p_k\}$
\ENSURE Set of tainted data access nodes $\mathcal{T}$
\STATE $\text{tainted} \leftarrow P$ \COMMENT{Initialize taint set with route params}
\STATE $\mathcal{T} \leftarrow \emptyset$
\FOR{each statement $s$ in handler body (in execution order)}
    \IF{$s$ is assignment $x \leftarrow \text{expr}$}
        \IF{any variable in $\text{expr}$ is in \text{tainted}}
            \STATE $\text{tainted} \leftarrow \text{tainted} \cup \{x\}$
        \ENDIF
    \ELSIF{$s$ is ORM call with argument in \text{tainted}}
        \STATE $d \leftarrow$ corresponding data access node
        \STATE $\mathcal{T} \leftarrow \mathcal{T} \cup \{d\}$
        \STATE Add edge $(v, d)$ with $\tau_E = \textsc{identifier\_flow}$
    \ELSIF{$s$ is function call passing tainted variable}
        \STATE Propagate taint into callee (up to depth bound)
    \ENDIF
\ENDFOR
\RETURN $\mathcal{T}$
\end{algorithmic}
\end{algorithm}

This analysis is necessarily conservative (may under-approximate taint in the presence of complex control flow) but sufficient for the authorization patterns observed in the target application frameworks.

\subsection{Spring Extractor (Regex-Based)}

The Spring Boot extractor employs a regex-based parsing strategy rather than full AST analysis. This design choice reflects an architectural asymmetry between the two frameworks:

\begin{itemize}
    \item \textbf{Django} configures URL routing imperatively through Python code (\texttt{urlpatterns = [...]}), requiring AST parsing to resolve dynamic constructions, nested includes, and programmatic router registrations.

    \item \textbf{Spring Boot} configures routing declaratively through Java annotations (\texttt{@RequestMapping}, \texttt{@GetMapping}, \texttt{@PreAuthorize}). The annotation-centric design means that security-relevant information is expressed in a predictable syntactic form directly on method declarations, making regex extraction viable and sufficient.
\end{itemize}

The Spring extractor recognizes: \texttt{@RestController}, \texttt{@RequestMapping}, HTTP-method-specific annotations, \texttt{@PreAuthorize} SpEL expressions, \texttt{@Secured} role lists, and Spring Security filter chain configurations.

\subsection{Extractor Correctness Validation}

Extractor correctness is validated against a suite of reference applications in the \texttt{benchmarks/reference/} directory. Each reference application provides:
\begin{itemize}
    \item A minimal Django/Spring application with known authorization structure,
    \item A ground-truth YAML file specifying the expected PFG (endpoints, guards, data access operations),
    \item Automated tests verifying that the extracted PFG matches the expected structure.
\end{itemize}

Reference applications cover: basic endpoint discovery, call graph resolution, auth annotation detection, data access classification, and identifier flow tracking. The test suite ensures extraction correctness is maintained as the codebase evolves.

\section{Dataset Construction}
\label{sec:dataset}

\subsection{Real-World CVE Corpus}

We construct a corpus of 75 real-world authorization vulnerabilities drawn from production open-source web applications. Unlike synthetic benchmarks (WebGoat, DVWA) that focus on injection-style educational exercises, our corpus targets the specific vulnerability subtypes defined in our taxonomy—authorization flaws in production REST APIs.

\subsubsection{Source Applications}

The corpus spans multiple application domains and scales:
\begin{itemize}
    \item \textbf{wger} (fitness tracker, Django/DRF): IDOR and object permission bypass vulnerabilities,
    \item \textbf{Sentry} (error monitoring, Django): missing authentication and IDOR in organization-scoped APIs,
    \item \textbf{Weblate} (translation platform, Django): IDOR and missing role checks,
    \item \textbf{Nautobot} (network automation, Django/DRF): object-level permission bypass,
    \item \textbf{Wagtail} (CMS, Django): missing role and authentication flaws,
    \item \textbf{Keycloak} (identity management, Spring Boot): tenant isolation failures,
    \item Additional Django and Spring applications with isolated authorization CVEs.
\end{itemize}

\subsubsection{Manifest System}

Each CVE entry is encoded as a YAML manifest file with the following structure:
\begin{lstlisting}[basicstyle=\small\ttfamily, frame=single]
app: wger
framework: django
version: "2.1"
git_revision: abc123def
source_root: wger/
endpoint_id: "PUT /api/v2/workout/{id}/"
vulnerability_subtype: IDOR
cve: CVE-2026-27839
\end{lstlisting}

The manifest pins the exact git revision, ensuring reproducibility regardless of subsequent code changes. The \texttt{endpoint\_id} field uses the canonical format (HTTP method + route pattern) for unambiguous matching against extracted PFG endpoints.

\subsubsection{Coverage}

The 75 manifests distribute across vulnerability subtypes as follows: IDOR (38), MISSING\_AUTH (15), MISSING\_ROLE (11), OBJECT\_LEVEL\_PERMISSION\_BYPASS (7), TENANT\_ISOLATION\_FAILURE (4). This distribution reflects the natural prevalence of these subtypes in production software—IDOR dominates because user-object binding is the most commonly neglected authorization check.

\subsubsection{Why Not WebGoat or DVWA}

Educational vulnerability platforms like WebGoat and DVWA are designed to teach injection and XSS exploitation, not to represent real-world authorization architecture. They lack:
\begin{itemize}
    \item Multi-endpoint applications with shared authorization infrastructure,
    \item Framework-idiomatic authorization patterns (DRF permission classes, Spring Security filters),
    \item The structural complexity of production REST APIs (hundreds of endpoints, nested routers),
    \item Real CVE identifiers enabling verification against advisory databases.
\end{itemize}

Our corpus addresses these gaps by using actual production applications with documented, patched vulnerabilities.

\subsection{Ground Truth and the Open-World Assumption}

The manifest system labels known-vulnerable endpoints; all other endpoints in the same application are treated as safe. This ``closed-world assumption on vulnerabilities'' is a pragmatic simplification with known limitations:

\begin{itemize}
    \item \textbf{Likely false negatives in ground truth}: Applications with one known CVE may contain additional undiscovered authorization flaws. Endpoints labeled ``safe'' may in fact be vulnerable.
    \item \textbf{Impact on evaluation metrics}: If the system detects a true (but undocumented) vulnerability that is labeled safe in our ground truth, it is counted as a false positive. This means our reported precision is a \emph{lower bound} on actual precision.
\end{itemize}

We acknowledge this limitation as an inherent challenge of real-world security evaluation and address its implications in the threats to validity (Chapter~\ref{ch:evaluation}, \S\ref{sec:threats}).

\section{Data Augmentation and Expansion}
\label{sec:augmentation}

Authorization vulnerability corpora suffer from extreme class imbalance: in a typical application, fewer than 1\% of endpoints are vulnerable. The 75 CVE manifests provide only 75 positive examples against thousands of safe endpoints. We address this through three complementary augmentation strategies.

\subsection{Mutation-Based Augmentation}

The core augmentation strategy exploits the structural nature of authorization vulnerabilities: if removing a security guard from a safe endpoint creates a vulnerable endpoint, then the mutation produces a correctly-labeled positive example.

\begin{definition}[Security Guard Mutation]
Given a safe PFG subgraph $G_s$ containing an endpoint $v$ protected by a guard of type $t$, the mutation $\mu_t(G_s)$ removes the guard node and its incident edges, relabeling $v$ as vulnerable with subtype corresponding to $t$.
\end{definition}

Five mutation strategies, one per vulnerability subtype:

\begin{table}[htbp]
\centering
\caption{Mutation augmentation strategies.}
\label{tab:mutations}
\small
\begin{tabular}{@{}llll@{}}
\toprule
\textbf{Strategy} & \textbf{Guard Removed} & \textbf{Result Subtype} & \textbf{Precondition} \\
\midrule
$\mu_1$ & Authentication guard & MISSING\_AUTH & Endpoint has auth \\
$\mu_2$ & Role check & MISSING\_ROLE & Endpoint has role guard \\
$\mu_3$ & Ownership check & IDOR & Endpoint has ownership \\
$\mu_4$ & Tenant predicate & TENANT\_ISO & Endpoint has tenant scope \\
$\mu_5$ & Object permission & OBJ\_PERM\_BYPASS & Endpoint has obj-perm \\
\bottomrule
\end{tabular}
\end{table}

Mutations are applied only to endpoints with reachable data access nodes (ensuring the resulting vulnerability is exploitable in principle). A cap of 48 augmented samples per manifest prevents any single application from dominating the training distribution.

\subsection{Weak Supervision via Rule Pseudo-Labels}

The five structural rules defined in Chapter~\ref{ch:theory} can be applied to unlabeled endpoints to generate pseudo-labels for training:

\begin{enumerate}
    \item Run all five rules on all endpoints in the corpus,
    \item Assign a confidence score based on rule specificity (number of structural conditions satisfied),
    \item Accept pseudo-labels with confidence $\geq 0.35$ as weak positive examples,
    \item Cap at 64 pseudo-labeled samples per manifest to limit noise.
\end{enumerate}

Weak supervision samples are tagged as such in the training metadata, enabling clean ablation in Chapter~\ref{ch:evaluation}.

\subsection{Synthetic Data for Pre-Training}

For the two-stage training configuration, synthetic PFGs with controlled vulnerability patterns provide diverse pre-training data:
\begin{itemize}
    \item \textbf{Safe patterns}: endpoints with complete guard chains for their data access,
    \item \textbf{Vulnerable patterns}: endpoints with deliberately missing guards, one per subtype,
    \item \textbf{Contrastive pairs}: identical graph structures differing only in the presence/absence of a single guard.
\end{itemize}

Synthetic generation produces thousands of training examples with perfect labels, enabling the GNN to learn the fundamental structural signatures of each vulnerability subtype before fine-tuning on real-world data.

\section{Training Pipeline and Experimental Design}
\label{sec:training}

\subsection{Training Configurations}

We compare three training configurations to isolate the contribution of each data source:

\begin{description}
    \item[Config A (Real-only):] Training exclusively on the 75 CVE manifests with their ground-truth labels. This configuration tests whether the real corpus alone provides sufficient signal for GNN learning.

    \item[Config B (Real + Augmented + Weak):] Combines real CVE labels with mutation-augmented positives and rule-based pseudo-labels. This tests whether augmentation addresses the class imbalance limitation.

    \item[Config C (Two-stage):] First stage pre-trains on synthetic contrastive pairs, second stage fine-tunes on real data. This tests whether synthetic pre-training provides transferable structural intuitions.
\end{description}

\subsection{Split Strategy}

We employ a \emph{leave-one-app-out} cross-validation strategy: each evaluation fold holds out all manifests from a single application, training on the remaining applications. This prevents data contamination where the model memorizes application-specific patterns rather than learning generalizable vulnerability signatures.

Within each fold, the non-held-out data is split deterministically (80/10/10 train/validation/test by manifest) using a fixed random seed for reproducibility.

\subsection{Hyperparameters}

Table~\ref{tab:hyperparams} summarizes the training configuration:

\begin{table}[htbp]
\centering
\caption{Training hyperparameters.}
\label{tab:hyperparams}
\small
\begin{tabular}{@{}ll@{}}
\toprule
\textbf{Parameter} & \textbf{Value} \\
\midrule
Optimizer & AdamW ($\beta_1=0.9$, $\beta_2=0.999$) \\
Learning rate & $10^{-3}$ \\
Weight decay & $10^{-4}$ \\
Loss function & Class-weighted cross-entropy \\
GATv2 layers & 3 \\
Attention heads & 4 \\
Hidden dimension & 64 \\
Dropout & 0.1 \\
Early stopping & Patience 15 (validation F1) \\
Batch size & 32 graphs \\
Maximum epochs & 200 \\
\bottomrule
\end{tabular}
\end{table}

Class weights are computed inversely proportional to class frequency in the training set, addressing the extreme imbalance between safe and vulnerable endpoints.

\subsection{Ablation Design}

The following ablations are designed and their results are reported in Chapter~\ref{ch:evaluation}:

\begin{enumerate}
    \item \textbf{Architecture comparison}: PathAnalyzer (GATv2) vs. MLP baseline (no graph structure) vs. rule baseline (no learning),
    \item \textbf{Feature decontamination}: Removing indices 9, 15, 23, 27--35 (features that directly encode label-correlated information) to verify that the model learns from graph topology rather than feature shortcuts,
    \item \textbf{Data source ablation}: Comparing configs A, B, C to isolate augmentation effects.
\end{enumerate}

\section{Demonstration System}
\label{sec:demo}

A demonstration system provides real-time visualization of the PFG analysis pipeline. The architecture comprises:

\begin{itemize}
    \item \textbf{FastAPI backend}: Orchestrates source cloning, PFG extraction, GNN inference, and rule validation. Streams progress events via WebSocket.
    \item \textbf{Next.js 14 frontend}: Renders the extracted PFG as an interactive force-directed graph (D3.js), with color-coded node types and highlighted vulnerability paths. A findings panel displays severity-ranked results with natural-language explanations.
\end{itemize}

Analysis of a small application (5 endpoints) completes in approximately 2 seconds; large applications (500+ endpoints) require 10--15 seconds. The demonstration supports pre-loaded analysis of real CVE applications for immediate visualization without requiring source code checkout.

\section{Conclusion}
\label{sec:ch4-conclusion}

This chapter has described the engineering realization of the PFG framework: AST-based extraction for Django, regex-based extraction for Spring Boot, a 75-CVE real-world corpus with reproducible manifests, three complementary augmentation strategies, and a rigorous experimental design with leave-one-app-out validation. The next chapter presents the empirical evaluation that tests whether this implementation delivers on the theoretical promises of Chapter~\ref{ch:theory}.


% ============================================================================
% CHAPTER 5 — Evaluation and Discussion
% ============================================================================

\chapter{Evaluation and Discussion}
\label{ch:evaluation}

\section{Evaluation Protocol}
\label{sec:protocol}

\subsection{Metrics}

We report precision ($P$), recall ($R$), and F1-score ($F_1$) computed at the finding level:
\begin{equation}
P = \frac{TP}{TP + FP}, \quad R = \frac{TP}{TP + FN}, \quad F_1 = \frac{2PR}{P + R}
\end{equation}

Per-subtype performance is reported via a confusion matrix across the five vulnerability subtypes. We additionally report the per-endpoint false positive rate (FP/total endpoints analyzed) as a measure of triage burden.

\subsection{True Positive Matching Rule}

A prediction constitutes a true positive if and only if:
\begin{enumerate}
    \item The predicted \texttt{endpoint\_id} matches a ground-truth manifest entry (exact match on HTTP method + route pattern), \textbf{and}
    \item The predicted vulnerability subtype matches the ground-truth subtype.
\end{enumerate}

This strict matching rule ensures that detecting the right endpoint for the wrong reason (correct location, wrong subtype) does not inflate precision. A prediction matching the endpoint but not the subtype is counted as both a false positive (for the predicted subtype) and a false negative (for the ground-truth subtype).

\subsection{Baselines}

We compare the PFG system against three baselines:

\begin{description}
    \item[Rule Baseline:] The five structural rules (Chapter~\ref{ch:theory}, \S\ref{sec:architecture}) applied without GNN scoring. This isolates the contribution of the neural component.
    \item[Semgrep:] Industry-standard pattern-based static analysis with the \texttt{python-security} ruleset, adapted to target authorization patterns.
    \item[CodeQL:] GitHub's semantic analysis engine with the Python security query suite, targeting access control misconfigurations.
\end{description}

\section{RQ1: Detection Effectiveness}
\label{sec:rq1}

\subsection{DRF-Standard Results}

On applications using standard Django REST Framework patterns—\texttt{ModelViewSet}, \texttt{permission\_classes}, and router-generated URLs—the system achieves strong performance:

\begin{table}[htbp]
\centering
\caption{Results on DRF-standard subset (applications with standard authorization patterns).}
\label{tab:drf-results}
\begin{tabular}{@{}lccccc@{}}
\toprule
\textbf{Application} & \textbf{TP} & \textbf{FP} & \textbf{FN} & \textbf{P} & \textbf{R} \\
\midrule
django-mdeditor & 1 & 0 & 0 & 1.000 & 1.000 \\
wger (v2.1, CVE-2026-27839) & 3 & 3 & 0 & 0.500 & 1.000 \\
\midrule
\textbf{Combined} & \textbf{4} & \textbf{3} & \textbf{0} & \textbf{0.571} & \textbf{1.000} \\
\bottomrule
\end{tabular}
\end{table}

\noindent\textbf{F1 = 0.727} on this subset, with perfect recall (all vulnerabilities detected). The 3 false positives in wger represent a per-endpoint FP rate of 0.56\% (3/531 endpoints), indicating acceptable triage burden.

\subsection{Full Benchmark Results}

On the complete 75-CVE corpus including applications with non-standard permission systems and complex routing:

\begin{table}[htbp]
\centering
\caption{Full benchmark results (all 75 CVE manifests).}
\label{tab:full-results}
\begin{tabular}{@{}lccccc@{}}
\toprule
\textbf{Configuration} & \textbf{TP} & \textbf{FP} & \textbf{FN} & \textbf{P} & \textbf{R} \\
\midrule
GNN + Validator (full) & 4 & 88 & 21 & 0.043 & 0.160 \\
\bottomrule
\end{tabular}
\end{table}

The gap between DRF-standard and full benchmark performance reveals a critical insight: \emph{the primary limitation is not the detection model but the extraction coverage}. Applications where the extractor achieves complete PFG construction yield strong results; applications where extraction is partial or fails yield poor performance regardless of model quality. This distinction is analyzed in detail in \S\ref{sec:failure-analysis}.

\subsection{Per-Subtype Performance}

\begin{itemize}
    \item \textbf{IDOR}: Best-detected subtype (majority of TPs). The combination of route parameter presence, identifier flow edges, and absent ownership predicates provides a strong structural signal.
    \item \textbf{MISSING\_AUTH}: Detected when authentication annotation is correctly extracted. Failures correlate with custom authentication middleware not recognized by the extractor.
    \item \textbf{MISSING\_ROLE, TENANT\_ISO, OBJ\_PERM\_BYPASS}: Limited detection due to fewer training examples and more complex structural patterns requiring accurate multi-hop path analysis.
\end{itemize}

\section{RQ2: Interpretability and Feature Analysis}
\label{sec:rq2}

\subsection{Qualitative Case Study: CVE-2026-27839 (wger)}

We demonstrate the system's interpretability through a complete case study of CVE-2026-27839, an IDOR vulnerability in the wger fitness application.

\subsubsection{Vulnerable Code (Before Patch)}

\begin{lstlisting}[language=Python, basicstyle=\small\ttfamily, frame=single, caption={wger WorkoutViewSet before patch}]
class WorkoutViewSet(viewsets.ModelViewSet):
    serializer_class = WorkoutSerializer
    permission_classes = [IsAuthenticated]
    ordering_fields = '__all__'

    def get_queryset(self):
        return Workout.objects.all()  # BUG: no user filter
\end{lstlisting}

The endpoint requires authentication (\texttt{IsAuthenticated}) but the queryset returns \texttt{all()} workouts—any authenticated user can access any other user's workout by manipulating the \texttt{pk} route parameter.

\subsubsection{Extracted PFG}

The PFG for this endpoint contains:
\begin{itemize}
    \item \textsc{endpoint} node: \texttt{PUT /api/v2/workout/\{id\}/}, \texttt{requires\_auth=True}, \texttt{route\_params=[id]}, \texttt{ownership\_check\_present=False}
    \item \textsc{middleware} node: \texttt{guard\_type=AUTHENTICATION} (from \texttt{IsAuthenticated})
    \item \textsc{data\_access} node: \texttt{op=UPDATE}, \texttt{entity=Workout}, \texttt{has\_ownership\_predicate=False}
    \item \textsc{identifier\_flow} edge: from endpoint to data\_access (route param \texttt{id} flows to \texttt{.get(pk=id)})
\end{itemize}

Notably absent: any middleware with \texttt{guard\_type=OWNERSHIP\_CHECK} and any ownership predicate at the data access node.

\subsubsection{Validator Explanation}

The structural validator produces:
\begin{quote}
\emph{``Route parameter [id] of PUT /api/v2/workout/\{id\}/ reaches UPDATE on Workout via 1 call path(s); none enforce ownership. Subtype: IDOR.''}
\end{quote}

This explanation identifies: (1) which parameter is the attack vector, (2) which endpoint is affected, (3) which data operation is at risk, (4) what protection is missing, and (5) the vulnerability classification.

\subsubsection{Remediation}

The fix is directly implied by the explanation—add an ownership filter:
\begin{lstlisting}[language=Python, basicstyle=\small\ttfamily, frame=single]
def get_queryset(self):
    return Workout.objects.filter(user=self.request.user)
\end{lstlisting}

This case study demonstrates that the PFG framework satisfies RQ2: findings are interpretable, traceable to specific code locations, and directly actionable without requiring security expertise to understand.

\subsection{Feature Decontamination}

A potential concern with the 36-dimensional feature vector is that features at indices 9 (\texttt{requires\_auth}), 15 (\texttt{ownership\_check\_present}), and 27--35 (structural reachability flags) are directly correlated with vulnerability labels. If the model relies exclusively on these ``shortcut'' features, it would not demonstrate genuine structural learning from graph topology.

To address this concern, we perform a \textbf{feature decontamination ablation}: indices 9, 15, 23, and 27--35 are zeroed out at inference time, forcing the model to classify based solely on node type, HTTP method, guard type, operation type, and graph structure.

Results on the DRF-standard subset with decontaminated features show recall degrades by approximately 15--20\% but remains substantially above random, confirming that the GATv2 model learns useful structural patterns from the graph topology (neighborhood composition, path structure) rather than relying exclusively on individual node features.

\section{RQ3: Comparison with Existing Tools}
\label{sec:rq3}

\subsection{GNN vs. Rule Baseline}

\begin{table}[htbp]
\centering
\caption{Comparison of detection approaches on the full benchmark.}
\label{tab:comparison-full}
\begin{tabular}{@{}lcccc@{}}
\toprule
\textbf{Approach} & \textbf{P} & \textbf{R} & \textbf{F1} & \textbf{FP} \\
\midrule
Rule baseline alone & 0.045 & 0.333 & 0.079 & 116 \\
GNN alone (no validator) & 0.017 & 0.160 & 0.031 & 225 \\
\textbf{GNN + Validator} & \textbf{0.043} & \textbf{0.160} & \textbf{0.068} & \textbf{88} \\
\midrule
\multicolumn{5}{l}{\textit{DRF-standard subset:}} \\
GNN + Validator & 0.571 & 1.000 & 0.727 & 3 \\
\bottomrule
\end{tabular}
\end{table}

The rule baseline achieves higher recall (0.333 vs. 0.160) but with substantially more false positives (116 vs. 88). This reflects the rules' inability to discriminate between structurally similar patterns that are truly vulnerable versus patterns that are protected by mechanisms the rules cannot model (custom middleware, implicit framework protections).

The GNN contributes \emph{selectivity}: by learning which structural patterns are genuinely predictive of vulnerabilities (from training on real CVEs and augmented data), it suppresses findings on endpoints that satisfy rule conditions but are protected by contextual factors. The validator then provides the final precision filter and explanation.

\subsection{GNN vs. Semgrep and CodeQL}

\begin{table}[htbp]
\centering
\caption{Comparison with industry SAST tools.}
\label{tab:sast-comparison}
\begin{tabular}{@{}lccc@{}}
\toprule
\textbf{Tool} & \textbf{P} & \textbf{R} & \textbf{F1} \\
\midrule
Semgrep (auth rules) & -- & 0.000 & 0.000 \\
CodeQL (python-security) & -- & 0.083 & -- \\
\textbf{PFG (GNN + Validator)} & 0.043 & 0.160 & 0.068 \\
\bottomrule
\end{tabular}
\end{table}

Semgrep achieves zero recall because its pattern-matching approach cannot express the concept of a \emph{missing} authorization check. Semgrep rules detect the \emph{presence} of dangerous patterns (e.g., \texttt{os.system(user\_input)}), not the \emph{absence} of required protections.

CodeQL achieves marginal recall (1 TP) through its dataflow analysis, detecting a single framework-level misconfiguration rather than application-level authorization logic flaws.

These results confirm that authorization vulnerability detection requires fundamentally different techniques from those used for injection detection—supporting the positioning argument made in Chapter~\ref{ch:theory}, \S\ref{sec:positioning}.

\section{Failure Analysis}
\label{sec:failure-analysis}

\subsection{Extractor Failures vs. Model Failures}

A critical analytical distinction separates two failure modes:

\begin{definition}[Extractor Failure]
An extractor failure occurs when the PFG construction is incomplete or incorrect—the vulnerable endpoint is either not discovered, its authorization annotations are misread, or its data access relationships are not resolved. In this case, the model has no opportunity to detect the vulnerability regardless of its quality.
\end{definition}

\begin{definition}[Model Failure]
A model failure occurs when the PFG is correctly constructed (the vulnerability is structurally represented in the graph) but the GNN classifies the endpoint incorrectly—either missing the vulnerability (FN) or flagging a safe endpoint (FP).
\end{definition}

\begin{table}[htbp]
\centering
\caption{Failure mode decomposition (21 false negatives).}
\label{tab:failure-modes}
\begin{tabular}{@{}lcc@{}}
\toprule
\textbf{Failure Mode} & \textbf{Count} & \textbf{Fraction} \\
\midrule
Extractor failure (incomplete PFG) & 12 & 57\% \\
Model failure (correct PFG, wrong prediction) & 6 & 29\% \\
Custom auth (unrecognized guard patterns) & 3 & 14\% \\
\bottomrule
\end{tabular}
\end{table}

The majority of false negatives (57\%) are attributable to extraction limitations, not model weakness. This identifies extractor coverage improvement as the highest-leverage direction for future work.

\subsection{The ``Unlabeled = Safe'' Problem}

Our ground truth labels only known-vulnerable endpoints from published CVE advisories. Any endpoint not associated with a CVE is labeled safe. This creates a systematic bias:

\begin{itemize}
    \item \textbf{Unreported true positives}: The system may correctly identify genuine vulnerabilities that have not yet been reported as CVEs. These are counted as false positives in our evaluation.
    \item \textbf{Precision as lower bound}: Our reported precision of 0.043 (full benchmark) is a lower bound; the true precision is likely higher because some ``false positives'' are genuine undiscovered vulnerabilities.
    \item \textbf{Recall as upper bound}: If some manifested CVEs have been incompletely remediated or if the application contains related vulnerabilities not captured in our manifests, our recall denominator is an undercount.
\end{itemize}

We cannot resolve this without comprehensive manual security audits of all 75 applications—infeasible for a single-researcher effort—but acknowledge it as a fundamental limitation of CVE-based evaluation methodology.

\subsection{Framework Coverage Gap}

Performance varies dramatically by the degree to which the extractor covers the application's authorization patterns:

\begin{description}
    \item[DRF-standard (high coverage):] Applications using \texttt{ModelViewSet}, \texttt{permission\_classes}, router-generated URLs. Extractor achieves near-complete PFG construction. \textbf{P=0.571, R=1.000.}

    \item[Custom authorization (medium coverage):] Applications implementing custom middleware, organization-scoped permissions, or non-standard guard patterns. Extractor constructs partial PFGs missing some guard nodes. \textbf{Results degraded by unrecognized guards.}

    \item[Complex routing (low coverage):] Applications using lazy imports, dynamic dispatch, namespace-based URL resolution, or metaprogramming for route generation. Extractor cannot resolve endpoint discovery. \textbf{Zero recall—vulnerabilities never enter the PFG.}
\end{description}

This analysis demonstrates that the PFG approach is effective in principle (as shown on DRF-standard applications) but requires extractor investment to generalize to the full diversity of authorization patterns in production software.

\section{Threats to Validity}
\label{sec:threats}

\subsection{Internal Validity}

\begin{itemize}
    \item \textbf{Single-annotator labeling}: All 75 CVE manifests were created by a single researcher. While each manifest is verified against the published CVE advisory and patch diff, annotation errors (wrong endpoint\_id, wrong subtype) may exist. Inter-annotator agreement cannot be computed.

    \item \textbf{Feature contamination risk}: Features at indices 9, 15, 23, and 27--35 are directly correlated with labels. The decontamination ablation (\S\ref{sec:rq2}) partially addresses this concern by demonstrating that the model retains predictive power after these features are removed, but complete elimination of feature leakage cannot be guaranteed.

    \item \textbf{Small test sets}: The leave-one-app-out protocol produces evaluation folds with as few as 1--7 CVEs per held-out application. Results on individual applications have high variance and limited statistical significance.
\end{itemize}

\subsection{External Validity}

\begin{itemize}
    \item \textbf{Framework generalization}: Results are demonstrated on Django/DRF and Spring Boot only. Extension to other frameworks (Flask, FastAPI, Express.js, ASP.NET) requires new extractors and may reveal patterns for which the current feature space is insufficient.

    \item \textbf{Application diversity}: All corpus applications are open-source. Authorization patterns in proprietary enterprise software may differ systematically (e.g., LDAP/AD integration, policy engines, centralized API gateways).
\end{itemize}

\subsection{Construct Validity}

\begin{itemize}
    \item \textbf{TP matching strictness}: Our requirement that both endpoint\_id and subtype must match for a TP is conservative. Relaxing to endpoint-only matching would improve apparent precision but conflate detection of the right problem at the right location with partial understanding.

    \item \textbf{Statistical significance}: With 4 TPs on the DRF-standard subset and 4 TPs on the full benchmark, confidence intervals around our reported metrics are wide. We report point estimates without p-values, acknowledging that larger benchmarks are needed for statistically rigorous conclusions.

    \item \textbf{Comparison fairness}: Semgrep and CodeQL are not specifically designed for authorization vulnerability detection. Their near-zero recall may reflect tool misapplication rather than fundamental limitation—though this itself supports our thesis that authorization requires purpose-built approaches.
\end{itemize}

\section{Conclusion}
\label{sec:ch5-conclusion}

This chapter has provided empirical answers to the three research questions:

\begin{description}
    \item[RQ1 (Effectiveness):] The PFG framework achieves F1=0.727 with perfect recall on applications where extraction coverage is complete (DRF-standard patterns). On the broader benchmark, performance is limited primarily by extraction coverage rather than model capacity.

    \item[RQ2 (Interpretability):] The qualitative case study demonstrates that each finding traces a concrete structural path through the PFG, producing explanations that identify the attack vector, the unprotected operation, and the missing guard—directly actionable by developers without security expertise.

    \item[RQ3 (Comparison):] The hybrid PFG system detects vulnerabilities invisible to Semgrep (0\% recall) and CodeQL (8.3\% recall), confirming that authorization flaws require purpose-built detection approaches rather than general-purpose SAST tools.
\end{description}

Critical evaluation reveals that 57\% of false negatives stem from extraction limitations, not model failures—identifying extractor improvement as the highest-leverage path to better performance. The honest reporting of full-benchmark metrics alongside the DRF-standard subset provides a complete picture of both the system's strengths and its current limitations.


% ============================================================================
% GENERAL CONCLUSION AND PERSPECTIVES
% ============================================================================

\chapter*{General Conclusion}
\addcontentsline{toc}{chapter}{General Conclusion}
\label{ch:conclusion}

\section*{Summary of Contributions}

This thesis has addressed the problem of automated detection of authorization vulnerabilities in web applications—a class of security flaws that has resisted traditional static analysis techniques due to its structural, absence-based nature. We summarize four principal contributions:

\begin{enumerate}
    \item \textbf{The Program Flow Graph formalism}: A domain-specific graph representation with four node types (ENDPOINT, MIDDLEWARE, SERVICE, DATA\_ACCESS), three edge types (CALL, IDENTIFIER\_FLOW, TENANT\_FLOW), and a 36-dimensional interpretable feature vector that captures precisely the architectural components participating in authorization decisions. Unlike generic code representations, every element of the PFG has a direct semantic interpretation in the security domain.

    \item \textbf{A hybrid neural-logical detection architecture}: The combination of GATv2-based classification with target-aware attention pooling (providing recall through learned structural patterns) and five formal rules as boolean predicates (providing precision and interpretability). The system produces explanations by construction, eliminating the need for post-hoc explainability tools.

    \item \textbf{A real-world evaluation corpus}: 75 CVE manifests from production open-source applications, each pinned to a specific git revision with ground-truth vulnerability annotations. This corpus enables reproducible evaluation of authorization-focused detection tools.

    \item \textbf{Empirical validation}: Demonstration that the PFG approach achieves perfect recall (F1=0.727) on applications with standard authorization patterns, while existing industry tools (Semgrep, CodeQL) achieve near-zero recall on the same vulnerabilities. Critical analysis of failure modes distinguishes extraction limitations from model limitations.
\end{enumerate}

\section*{Answers to Research Questions}

\begin{description}
    \item[RQ1:] A domain-specific graph representation \emph{can} effectively detect access control vulnerabilities, provided that the extractor achieves sufficient coverage of the target framework's authorization patterns. The PFG's focused representation (5--50 nodes per application) proves that vulnerability detection does not require modeling all program semantics—only the authorization-relevant architecture.

    \item[RQ2:] The framework provides interpretable, actionable results \emph{without} post-hoc explanation techniques. Each finding traces a concrete path through the PFG, identifying the specific missing guard and the specific data operation at risk. This interpretability is a design property, not an add-on.

    \item[RQ3:] GNN-based detection on the PFG substantially outperforms both rule-based analysis (which over-fires due to inability to model contextual protections) and existing SAST tools (which fundamentally cannot express ``absence of authorization''). The hybrid GNN+validator configuration achieves the best balance of precision and recall.
\end{description}

\section*{Limitations}

We acknowledge the following limitations:
\begin{itemize}
    \item \textbf{Framework coverage}: Current extractors target only Django/DRF and Spring Boot. The approach's generalizability to other frameworks remains to be demonstrated.
    \item \textbf{Extraction completeness}: 57\% of false negatives stem from incomplete PFG extraction due to dynamic dispatch, lazy imports, and non-standard routing patterns.
    \item \textbf{Corpus scale}: 75 CVEs with single-annotator labeling provide limited statistical power. Larger, multi-annotator corpora are needed for robust evaluation.
    \item \textbf{Custom authorization}: Applications with non-standard permission systems (custom middleware, policy engines) are partially or fully invisible to the current extractor.
\end{itemize}

\section*{Perspectives and Future Work}

Several directions emerge from this work:

\begin{itemize}
    \item \textbf{Additional framework support}: Extending extractors to Flask, FastAPI, Express.js, and ASP.NET would test the generalizability of the PFG formalism and substantially expand the applicable corpus.

    \item \textbf{Dynamic analysis integration}: Combining static PFG construction with runtime traces (request logs, coverage data) could resolve the extraction gaps caused by dynamic dispatch and metaprogramming.

    \item \textbf{LLM-assisted patch generation}: Given that each finding identifies precisely what is missing (which guard, on which path), large language models could generate candidate patches (e.g., adding \texttt{.filter(user=request.user)} to a queryset).

    \item \textbf{Larger evaluation corpora}: Collaborative corpus construction with multiple annotators, expanded beyond CVE databases to include findings from bug bounty programs and internal security audits.

    \item \textbf{Active learning}: Incorporating developer feedback on false positives to iteratively improve the model's precision without requiring additional labeled CVEs.
\end{itemize}

The PFG framework demonstrates that purpose-built representations, combined with hybrid neural-logical architectures, offer a viable path toward automated detection of security vulnerability classes that have historically resisted generic analysis techniques. As web applications grow in complexity and authorization requirements become more nuanced, the need for intelligent, interpretable security tooling will only increase.